\documentclass[11pt, a4paper]{article}

\usepackage[margin=2.5cm]{geometry}
\usepackage[T1]{fontenc}
\usepackage[utf8]{inputenc}
\usepackage{palatino}
\usepackage{mathpazo}
\usepackage{microtype}
\usepackage{booktabs}
\usepackage{longtable}
\usepackage{array}
\usepackage{hyperref}
\usepackage{amsmath}
\usepackage{amssymb}
\usepackage{enumitem}
\usepackage{xcolor}
\usepackage{fancyhdr}
\usepackage{titlesec}
\usepackage{tcolorbox}
\usepackage{multirow}
\usepackage{caption}
\usepackage{mdframed}

% ── Colour palette (muted, academic) ────────────────────────────────────────
\definecolor{navy}{RGB}{18, 34, 66}
\definecolor{steel}{RGB}{55, 90, 135}
\definecolor{slategray}{RGB}{100, 120, 145}
\definecolor{palegray}{RGB}{243, 245, 248}
\definecolor{rulegray}{RGB}{160, 180, 200}
\definecolor{defblue}{RGB}{230, 236, 245}

% ── Hyperref ────────────────────────────────────────────────────────────────
\hypersetup{
  colorlinks=true,
  linkcolor=steel,
  urlcolor=steel,
  citecolor=steel,
  pdftitle={Pre-Registration Package: Testable Predictions from the Salience-Centered Paradigm},
  pdfauthor={Flyxion},
  pdfsubject={Pre-registration; Signal Processing; Auditory Perception; Associative Memory}
}

% ── Section formatting ───────────────────────────────────────────────────────
\titleformat{\section}
  {\large\bfseries\color{navy}}{}{0em}{}
  [\vspace{0.1em}{\color{rulegray}\hrule height 0.4pt}]
\titleformat{\subsection}{\normalsize\bfseries\color{steel}}{}{0em}{}
\titleformat{\subsubsection}{\normalsize\itshape\color{navy}}{}{0em}{}

% ── Header / footer ─────────────────────────────────────────────────────────
\pagestyle{fancy}
\fancyhf{}
\fancyhead[L]{\small\color{slategray} Pre-Registration Package}
\fancyhead[R]{\small\color{slategray} Salience-Centered Paradigm — Levels I–III}
\fancyfoot[C]{\small\color{slategray}\thepage}
\renewcommand{\headrulewidth}{0.4pt}
\renewcommand{\headrule}{\hbox to\headwidth{\color{rulegray}\leaders\hrule height \headrulewidth\hfill}}
\setlength{\headheight}{14pt}

% ── Box environments ─────────────────────────────────────────────────────────
\tcbuselibrary{skins, breakable}

% Prediction template box
\newtcolorbox{predbox}[1]{
  enhanced, breakable,
  title=#1,
  fonttitle=\bfseries\color{white}\small,
  colbacktitle=navy,
  colback=palegray,
  colframe=rulegray,
  boxrule=0.5pt,
  arc=2pt,
  top=3pt, bottom=3pt, left=6pt, right=6pt,
  titlerule=0pt,
  attach boxed title to top left={yshift=-2mm, xshift=5mm},
  boxed title style={colback=navy, arc=2pt, boxrule=0pt}
}

% Definition box
\newtcolorbox{defbox}{
  enhanced, breakable,
  colback=defblue,
  colframe=steel!60,
  boxrule=0.4pt,
  arc=2pt,
  left=6pt, right=6pt, top=3pt, bottom=3pt,
  fontupper=\small
}

% Note box
\newtcolorbox{notebox}{
  enhanced, breakable,
  colback=white,
  colframe=rulegray,
  boxrule=0.4pt,
  arc=2pt,
  left=6pt, right=6pt, top=3pt, bottom=3pt,
  fontupper=\small\itshape
}

% ── List settings ────────────────────────────────────────────────────────────
\setlist[itemize]{itemsep=1pt, topsep=3pt, leftmargin=1.5em}
\setlist[enumerate]{itemsep=1pt, topsep=3pt, leftmargin=1.5em}

% ── Paragraph spacing ────────────────────────────────────────────────────────
\setlength{\parindent}{14pt}
\setlength{\parskip}{0.4ex}

% ── Section spacing ──────────────────────────────────────────────────────────
\titlespacing*{\section}{0pt}{1.4ex plus 0.4ex minus 0.2ex}{0.8ex}
\titlespacing*{\subsection}{0pt}{1.0ex plus 0.3ex minus 0.2ex}{0.4ex}
\titlespacing*{\subsubsection}{0pt}{0.8ex plus 0.2ex minus 0.1ex}{0.2ex}

% ── Convenience commands ─────────────────────────────────────────────────────
\newcommand{\mem}{MEM$|$8}
\newcommand{\hcf}{HCF1}
\newcommand{\scp}{Salience-Centered Paradigm}
\newcommand{\dprime}{$d'$}

% ============================================================================
\begin{document}
\raggedbottom

% ── Title page ───────────────────────────────────────────────────────────────
\begin{titlepage}
\vspace*{1.0cm}
{\color{navy}\rule{\linewidth}{1.5pt}}
\vspace{0.4cm}

\begin{center}
{\LARGE\bfseries\color{navy}
Pre-Registration Package:\\[0.45em]
Testable Predictions Derived from the\\[0.45em]
Salience-Centered Paradigm (Levels~I–III)
}

\vspace{0.6cm}
{\large\color{steel}
Signal Processing $\cdot$ Auditory Perception $\cdot$ Associative Memory Architecture
}

\vspace{1.0cm}
\begin{tabular}{rl}
\textbf{Author:}       & Flyxion (independent researcher) \\[4pt]
\textbf{Components:}   & Marine Salience Algorithm; \hcf{} Compression; \mem{} \\[4pt]
\textbf{Scope:}        & Six falsifiable predictions (Sets A, B, C) \\[4pt]
\textbf{Suitable for:} & OSF, AsPredicted, Registered Reports tracks \\[4pt]
\textbf{Date:}         & May 2026
\end{tabular}

\vspace{1.0cm}
{\color{navy}\rule{\linewidth}{0.7pt}}
\vspace{0.4cm}

\begin{notebox}
\textbf{Scope of this document.}
This package concerns only the empirically tractable claims within Levels~I–III
of the \scp{} (signal processing, compression architecture, associative memory).
Level~IV claims (consciousness, authenticity loss, and related constructs) are not
pre-registered here; they require prior operationalization before experimental
design is possible.  The goal of this document is not to validate the broader
paradigm but to derive risky, testable predictions from it — predictions whose
failure would carry genuine evidential weight.
\end{notebox}

\end{center}
\vfill
{\color{navy}\rule{\linewidth}{1.5pt}}
\end{titlepage}

% ── Abstract ─────────────────────────────────────────────────────────────────
\thispagestyle{plain}
{\centering\large\bfseries\color{navy} Abstract\par}
\smallskip
\noindent
The \scp{} is a research programme proposing that temporal coherence — specifically,
structured micro-temporal variation (jitter) in harmonic oscillators — constitutes
signal-level information rather than noise, and that this insight motivates novel
approaches to perceptual salience estimation, audio compression, and phase-based
associative memory.  The paradigm spans four levels of increasing abstraction:
temporal signal processing (Level~I), compression architecture (Level~II),
memory architecture (Level~III), and speculative claims about consciousness and
authenticity (Level~IV).

This pre-registration package concerns exclusively Levels~I–III, which generate
six falsifiable predictions amenable to standard experimental and computational
methods.  For Levels~I and~II, predictions address (A1)~perceptual detection of
temporal jitter in harmonic complexes under spectrally controlled conditions,
(A2)~incremental predictive power of a coherence metric in auditory salience ratings,
(B1)~linear predictability of upper harmonic envelopes from the fundamental in
acoustic instruments, and (B2)~bitrate efficiency and perceptual transparency of
the \hcf{} compression scheme.  For Level~III, predictions address (C1)~quadratic
storage scaling in phase-encoded associative memory relative to a Hopfield baseline,
and (C2)~resistance to catastrophic forgetting under sequential learning.

Each template specifies a dependent variable, null hypothesis, statistical criterion,
sample size with power justification, analysis pipeline, and
confirmatory/exploratory distinction.  The package follows open-science norms for
pre-registration and is intended to be filed at OSF, AsPredicted, or an equivalent
repository.  Negative results are treated as substantive evidence; falsification of
individual predictions is explicitly distinguished from falsification of the paradigm
as a whole.

\noindent\textit{Keywords:} auditory salience; temporal coherence; harmonic compression;
associative memory; pre-registration; falsifiability; signal detection theory;
open science.

\newpage

% ── TOC ──────────────────────────────────────────────────────────────────────
\tableofcontents
\newpage

% ── 1. NOTATION AND TERMINOLOGY ───────────────────────────────────────────────
\section{Notation and Terminology}

The following definitions are fixed for the entirety of this document.
Terminological consistency matters for pre-registrations because post-hoc
reinterpretation of key terms can obscure whether a prediction was confirmed or
disconfirmed.

\begin{defbox}
\textbf{Temporal jitter.}  Cycle-to-cycle variation in the period of a quasi-periodic
waveform, quantified as $J = \sigma(T) / \mu(T)$, where $T$ denotes the set of
measured inter-peak intervals over a specified time window (default: 200~ms).
Jitter is expressed as a percentage of mean period.

\smallskip
\textbf{Marine coherence metric ($\kappa$).}  A normalised inverse-jitter quantity
defined as $\kappa = 1 / (J + \varepsilon)$, where $\varepsilon > 0$ is a small
regularising constant (default: $10^{-3}$).  Higher $\kappa$ indicates more
temporally stable oscillation.  The ``Marine'' qualifier refers to the salience
algorithm in which this metric originates; it does not carry implications about
consciousness or authenticity.

\smallskip
\textbf{Spectral prominence ($P$).}  Level of a partial above the local noise
floor, measured in dB: $P = L_\text{partial} - L_\text{noise}$.

\smallskip
\textbf{HCF1 (Harmonic Conformity Flag, 1-bit).}  A compression scheme storing a
master amplitude envelope for the fundamental ($e_1$) and, for each upper harmonic
$h = 2, \ldots, H$, a single bit indicating conformity to a pre-calibrated linear
prediction $\hat{e}_h = a_h e_1 + b_h$.  Non-conforming frames store an explicit
residual.

\smallskip
\textbf{\mem{} (Modulated Encoding Memory, 8 wave functions per node).}  A proposed
associative memory architecture in which pairwise associations are encoded as phase
relationships among complex sinusoidal wave functions.  The hypothesis is that $N$
wave functions can store $O(N^2)$ associations, contrasting with the $O(N)$
capacity of standard Hopfield networks.

\smallskip
\textbf{Constraint horizon.}  The requirement that any empirical claim specify in
advance what evidence would count against it.  A framework that can accommodate
every possible outcome is not a scientific theory; it is a descriptive vocabulary.
All six predictions in this package are designed to have a clear constraint horizon.

\smallskip
\textbf{Retention (\%).}  In sequential learning experiments (C2): the ratio of
recall accuracy after learning a second association set to recall accuracy before,
expressed as a percentage.

\smallskip
\textbf{Confirmatory vs.\ exploratory.}  Analyses declared in this document before
data collection are \emph{confirmatory}.  Analyses conducted after data collection,
whether planned or opportunistic, are \emph{exploratory} and must be labelled as
such in any resulting publication.
\end{defbox}

\bigskip

\noindent Throughout this document, notation is standardised as follows:

\noindent\begin{minipage}{\textwidth}
\centering
\renewcommand{\arraystretch}{1.3}
\begin{tabular}{ll}
\toprule
\textbf{Symbol / term} & \textbf{Meaning} \\
\midrule
$f_0$              & Fundamental frequency \\
$h$                & Harmonic index ($h = 1$ = fundamental) \\
$e_h(t)$           & Amplitude envelope of harmonic $h$ at time $t$ \\
$J$                & Jitter (normalised period SD) \\
$\kappa$           & Marine coherence metric \\
$P$                & Spectral prominence (dB) \\
$d'$               & Sensitivity index (signal detection theory) \\
$R^2$              & Coefficient of determination \\
$\Delta R^2$       & Incremental explained variance \\
$N$                & Number of wave functions in \mem{} \\
$M$                & Number of items forming $M^2$ pairwise associations \\
\mem{}             & Modulated Encoding Memory architecture \\
\hcf{}             & Harmonic Conformity Flag (1-bit) compression scheme \\
\bottomrule
\end{tabular}
\end{minipage}

% ── 2. SCIENTIFIC POSITIONING ─────────────────────────────────────────────────
\section{Scientific Positioning}

The predictions in this package sit within three established research traditions.
Anchoring them explicitly prevents the framework's novel terminology from appearing
to arise in a vacuum.

\subsection{Prediction Set A — Psychoacoustics and Temporal Coding}

The question of whether fine temporal structure contributes to auditory perception
beyond what is captured by the spectral envelope has a long history.  Phase locking
in the auditory nerve supports temporal coding up to approximately 4–5~kHz (Rose
et al., 1967; Palmer \& Russell, 1986).  The role of micro-temporal variation in
the perception of naturalness and timbral identity has been studied in the context
of vibrato, mechanical fluctuations in musical instruments, and the perceptual
consequences of phase randomisation (McAdams, 1984; Beauchamp, 1993).  Auditory
salience models typically operationalise salience in terms of spectral energy,
spectral flux, or onset strength (Kayser et al., 2005; Klapuri, 2006); the
Marine coherence metric represents a proposal that temporal stability constitutes
an orthogonal salience dimension not captured by these spectral proxies.
Prediction~A2 tests this directly as an incremental regression effect.

\subsection{Prediction Set B — Perceptual Audio Coding and Harmonic Modelling}

Harmonic signal models underlie a broad family of audio codecs and analysis
methods, including sinusoidal modelling (McAulay \& Quatieri, 1986), source-filter
approaches to speech coding (Rabiner \& Schafer, 1978), and structured prediction
in transform codecs.  The observation that harmonic amplitude envelopes co-vary
across a note (Prediction~B1) is related to but distinct from standard spectral
envelope coding: HCF1 exploits inter-harmonic temporal correlation rather than
across-frequency spectral shape.  In this sense it is closer in spirit to
predictive coding of envelope trajectories (Spanias, 1994) than to masked-band
discard strategies used in perceptual codecs such as MP3 and AAC.

\subsection{Prediction Set C — Associative Memory and Attractor Networks}

Hopfield networks (Hopfield, 1982) are the canonical model of content-addressable
memory with a known storage capacity of approximately $0.138N$ patterns for $N$
binary units (Amit et al., 1985).  Complex-valued and phase-coded extensions of
Hopfield-type networks have been studied theoretically (Noest, 1988; Jankowski
et al., 1996; Aoyagi, 1995) and motivate the claim that phase relationships among
oscillatory units can encode pairwise associations more efficiently than binary
state vectors.  The MEM$|$8 architecture is best understood as a proposal within
this family of complex-valued associative networks.  Predictions C1 and~C2 are
designed to be interpretable within this existing literature: the quadratic
capacity claim (C1) is a strong quantitative prediction relative to known Hopfield
scaling, and the catastrophic forgetting result (C2) connects to the well-studied
stability-plasticity dilemma (Grossberg, 1987; McCloskey \& Cohen, 1989).

% ── 3. CROSS-LEVEL ARCHITECTURAL UNITY ────────────────────────────────────────
\section{Cross-Level Architectural Unity}

Although Prediction Sets A, B, and C address distinct experimental domains, they
are unified by a single underlying principle: \emph{structured redundancy under
physical constraint}.

Level~I (Prediction Set A) concerns the \textbf{extraction} of structurally
informative temporal patterns from harmonic signals.  The claim is that jitter is
not noise to be discarded but signal to be retained because it encodes the physical
coupling state of the oscillating system.

Level~II (Prediction Set B) concerns the \textbf{representation} of correlated
harmonic structure.  Given that upper harmonics are driven by shared physical
dynamics, their amplitude envelopes are expected to be highly predictable from the
fundamental.  \hcf{} exploits this predictability to achieve compact, lossless
(within-frame) description of harmonic temporal evolution.

Level~III (Prediction Set C) concerns \textbf{retrieval} under structured relational
encoding.  If associative memory is implemented via phase relationships among
oscillatory units, then the information-theoretic capacity of the representation
scales differently from binary state vectors — specifically, quadratically with the
number of wave functions.

All three levels therefore concern constrained representation under physically or
structurally motivated redundancy.  The framework's contribution, if confirmed,
would be identifying that this principle applies coherently across the domains of
perception (Level~I), compression (Level~II), and memory (Level~III).

It is important to note that the three levels are architecturally related but
\emph{empirically independent}.  Falsification of one prediction set does not
imply falsification of the others.  See Section~\ref{sec:negative} for a systematic
treatment of what each negative result would and would not entail.

% ── 4. EXPLAINABILITY AS AN INDEPENDENT CONTRIBUTION ─────────────────────────
\section{Explainability as an Independent Contribution}

A common implicit assumption in systems like \mem{} is that the primary measure
of success is capacity or efficiency relative to existing architectures.  This
framing sets a high bar: if the quadratic capacity claim in C1 is not confirmed,
the architecture might appear to offer nothing.  This section records a secondary
contribution claim that does \emph{not} depend on the capacity results.

Phase-based associative memory encodes relationships as geometric objects in the
complex plane.  Each stored association $(i, j)$ corresponds to an angular offset
$\Delta\phi_{ij}$ between two wave functions.  This has two interpretive
consequences that standard binary Hopfield networks do not share:

\begin{enumerate}
  \item \textbf{Geometric interpretability.}  The strength of an association can
    be read directly from the phase offset magnitude; the direction of retrieval
    (from $i$ to $j$ vs.\ $j$ to $i$) corresponds to the sign of the angle.
    A trained weight matrix in a binary Hopfield network offers no analogous
    geometric decomposition without additional analysis.

  \item \textbf{Graded similarity.}  Items with similar phase profiles are
    geometrically close in the complex plane, enabling a natural similarity metric
    without additional distance computation.  This may be useful for approximate
    nearest-neighbour retrieval in applications where exact matches are rare.
\end{enumerate}

These properties hold regardless of whether the capacity results in C1 are
confirmed.  If C1 fails — if the quadratic scaling does not materialize — the
architecture still offers a transparent, geometrically interpretable
representation of associative structure that may be preferable to binary
attractor networks in settings where interpretability matters (e.g., embedded
systems, explainable AI pipelines, or educational demonstrations of associative
memory).

This claim is explicitly \emph{not} pre-registered as a falsifiable prediction
in the present document, because interpretability is a design property, not an
experimental outcome.  It is recorded here so that the broader value of the
\mem{} architecture is not contingent on the specific capacity thresholds in C1.

% ── 5. OPERATIONAL BOUNDARY CONDITIONS ────────────────────────────────────────
\section{Operational Boundary Conditions}
\label{sec:boundaries}

A well-bounded framework is more tractable, not weaker.  The following are explicit
non-claims of the Salience-Centered Paradigm as operationalized in this package.

\begin{enumerate}[label=\textbf{B\arabic*.}, leftmargin=2.2em]

  \item \textbf{No claim to proof of consciousness.}  Nothing in Levels~I–III
    entails or supports claims about conscious experience, phenomenal authenticity,
    or subjective perception of ``cognitive field'' effects.  The Marine coherence
    metric is a signal-processing quantity.

  \item \textbf{No rejection of standard psychoacoustic masking.}  The framework
    does not claim that psychoacoustic masking models (as used in MP3, AAC, and
    related codecs) are incorrect.  Prediction~A1 is compatible with masking-based
    models; it tests whether jitter carries additional perceptual information
    \emph{beyond} masked-band predictions.

  \item \textbf{No claim of current engineering superiority.}  \hcf{} compression
    is not claimed to outperform existing perceptual codecs in general.
    Prediction~B2 tests a specific, narrow claim about bitrate reduction for
    harmonic envelope coding, not end-to-end audio quality.

  \item \textbf{No claim of biological plausibility for \mem{}.}  The \mem{}
    architecture is a mathematical model.  Predictions C1 and C2 are claims about
    a defined computational system, not about neural implementation in biological
    tissue.

  \item \textbf{No claim that Level~IV is false.}  This document takes no position
    on whether Level~IV claims about consciousness and authenticity are correct.  They
    are excluded solely because they lack sufficient operationalization to support
    pre-registered experimental design at this time.

  \item \textbf{No claim that confirmation of Levels~I–III entails Level~IV.}
    Even if all six predictions are confirmed, that result would not constitute
    evidence for the Level~IV speculative claims.  The inferential boundary is sharp.

\end{enumerate}

% ── 5. STATISTICAL PHILOSOPHY AND CONFIRMATORY LOGIC ─────────────────────────
\section{Statistical Philosophy and Confirmatory Logic}

\subsection{Confirmatory versus exploratory analyses}

This package follows the \emph{strong} form of pre-registration: all confirmatory
analyses are fully specified prior to any data collection, including dependent
variables, null hypotheses, statistical tests, alpha levels, sample sizes, and
exclusion criteria.  Analyses not specified here are \emph{exploratory} by default
and must be labelled as such.  Exploratory analyses are scientifically valuable and
are not discouraged; the distinction matters solely for inferential status.

\subsection{Effect thresholds and estimation}

Null hypothesis tests are supplemented by explicit effect thresholds wherever
possible (e.g., $d' > 0.5$, $\Delta R^2 \geq 0.15$, retention $\geq 90\%$).
These thresholds are motivated by theoretical and practical significance, not by
statistical convenience.  The goal is to ask not only whether an effect exists but
whether it is large enough to matter for the claims of the paradigm.

\subsection{Signal detection theory (Predictions A1, A2)}

Signal detection theory (SDT) provides the appropriate framework for perceptual
discrimination studies because it dissociates sensitivity ($d'$) from response bias
($c$).  A listener who is more willing to report jitter (liberal criterion) will
produce more hits but also more false alarms; $d'$ is unaffected by this bias.
This is important for Prediction~A1 because motivation, attention, and individual
criterion placement could otherwise confound a simple proportion-correct measure.

\subsection{Incremental regression (Prediction A2)}

The decision to test $\Delta R^2$ rather than a raw correlation reflects the
incremental logic of the Marine coherence hypothesis: the claim is that coherence
predicts salience \emph{beyond} spectral prominence, not that it predicts salience
in isolation.  A model comparison that does not control for prominence would not
test the paradigm's actual claim.

\subsection{Replication priority}

Narrative coherence of the paradigm is explicitly subordinated to replication
priority.  If a prediction fails, the result stands.  Post-hoc reinterpretation
of failed predictions as consistent with the paradigm is not permitted within the
confirmatory framework; it must be declared exploratory.

% ── 6. RELATION TO EXISTING COMPRESSION PARADIGMS ────────────────────────────
\section{Relation to Existing Compression Paradigms}

The \hcf{} scheme (Prediction Set B) is most clearly understood by contrast with
existing approaches to audio compression.

\noindent\begin{minipage}{\textwidth}
\centering
\renewcommand{\arraystretch}{1.4}
\begin{tabular}{p{3.2cm} p{4.2cm} p{4.8cm}}
\toprule
\textbf{Paradigm} & \textbf{What is discarded / predicted} & \textbf{Contrast with HCF1} \\
\midrule
Transform coding (MP3, AAC) &
  Masked spectral bands; fine temporal structure &
  HCF1 retains all temporal structure; no masking model used \\
Linear predictive coding (LPC) &
  Spectral envelope; residual coded separately &
  HCF1 operates on amplitude envelopes, not the waveform directly \\
Sinusoidal modelling &
  Non-sinusoidal residual; parameters tracked per partial &
  HCF1 exploits inter-harmonic envelope correlation, not per-partial tracking \\
Predictive envelope coding &
  Temporal redundancy across frames &
  HCF1 exploits redundancy \emph{across harmonics} within a frame \\
\bottomrule
\end{tabular}
\captionof{table}{Conceptual comparison of \hcf{} with established compression paradigms.}
\end{minipage}
The distinguishing feature of \hcf{} is the hypothesis that upper harmonic
amplitude envelopes are predictable from the fundamental with sufficient accuracy
to permit 1-bit flag encoding.  This is an empirical claim (Prediction~B1) that
must be confirmed before the compression scheme can be evaluated (Prediction~B2).
The two predictions are therefore logically ordered: B1 is the necessary prior
condition for B2.

% ── 7. COMPUTATIONAL COMPLEXITY OF MEM|8 ─────────────────────────────────────
\section{Computational Complexity and Scalability of \mem{}}

Prediction Set C involves a computational architecture that may be unfamiliar to
readers outside the associative memory literature.  The following summary clarifies
the relevant asymptotic properties and their experimental implications.

\subsection{Storage complexity}

A \mem{} system with $N$ wave functions has two distinct storage components.
First, the $N$ oscillator states $\{\phi_k\} \in \mathbb{C}^N$ contribute $O(N)$
degrees of freedom.  Second, an antisymmetric pairwise offset matrix
$\Delta_{ij} \in \mathbb{R}^{N \times N}$ (zero diagonal) contributes
$O(N^2)$ entries — one target angular offset per ordered pair.  It is this
matrix, not the oscillators alone, that provides the hypothesised relational
capacity.  The claim that $O(N^2)$ associations can be stored is therefore
a claim about constraint-satisfaction in the joint oscillator--offset system,
not a claim that phase continuity alone implies quadratic capacity.

This distinction matters for interpretation: if C1 is confirmed, the result
would suggest that a phase-relational constraint system exploits pairwise
angular structure efficiently and exhibits sharp retrieval transitions under
interference.  If C1 fails, the failure would most likely indicate that the
relaxation algorithm cannot simultaneously satisfy $O(N^2)$ constraints given
only $O(N)$ oscillator degrees of freedom.

For comparison, the standard Hopfield network stores at most $\approx 0.138N$
patterns reliably before interference dominates (Amit et al., 1985).

\subsection{Retrieval complexity}

Retrieval in \mem{} involves nearest-neighbour search in phase space.  For a
full pairwise comparison this is $O(N^2)$, though approximate retrieval schemes
could reduce this.  Predictions~C1 and~C2 are agnostic to the retrieval algorithm;
they concern accuracy at convergence, not speed.

\subsection{Robustness under noise}

A known limitation of phase-coded associative memories is sensitivity to phase
noise: small errors in the representation of wave functions can produce retrieval
failures.  The simulations in Predictions~C1 and~C2 use floating-point arithmetic
and therefore understate real-world noise sensitivity.  Any practical deployment
would require additional robustness analysis.  The predictions in this package are
therefore best understood as establishing \emph{upper bounds} on capacity and
retention under idealised conditions.

\subsection{Scalability}

Experiments~C1 and~C2 use $N \in \{64, 128, 256\}$.  These are toy-scale models.
Whether the claimed advantages persist at $N \gtrsim 10^4$ (required for practical
application) is an open question that lies beyond the scope of this pre-registration.

% ── 8. PREAMBLE AND FRAMEWORK ─────────────────────────────────────────────────
\section{Preamble and Epistemic Framework}

The predictions collected in this package are derived from the technically grounded
strata of the \scp{}.  The organizing principle throughout is the
\emph{constraint horizon}: a well-formed empirical claim must specify in advance
what evidence would count against it.  The framework generates these predictions
not as confirmatory demonstrations but as genuine empirical risks.  A prediction
set whose failure would be shrugged off is not a scientific prediction; each of the
six templates below is designed so that a negative outcome would carry clear
evidential weight against the corresponding level of the paradigm.

\noindent\begin{minipage}{\textwidth}
\centering
\renewcommand{\arraystretch}{1.4}
\begin{tabular}{p{1.3cm} p{4.8cm} p{5.5cm}}
\toprule
\textbf{Level} & \textbf{Core claim} & \textbf{Status} \\
\midrule
I   & Temporal jitter encodes structural information beyond masking &
      Testable; Prediction Set A \\
II  & Harmonic envelopes are linearly predictable from the fundamental;
      1-bit flag compression is efficient &
      Testable; Prediction Set B \\
III & Phase-based associative memory yields quadratic capacity and resists
      catastrophic forgetting &
      Testable; Prediction Set C \\
IV  & Ultrasonic frequencies alter brain coherence; spectral fracture produces
      ``authenticity loss'' &
      \emph{Excluded — requires prior operationalization} \\
\bottomrule
\end{tabular}
\end{minipage}

\begin{notebox}
\textbf{Data availability.}  All stimuli, code, and anonymized data will be
deposited to an open repository (OSF, Zenodo, or equivalent) upon study completion.
Deviations from the pre-registered analysis plan must be declared in print.
\end{notebox}

\begin{notebox}
\textbf{Conflicts of interest.}  The author has a theoretical investment in the
\scp{}.  Independent replication is therefore especially encouraged.
\end{notebox}

% ── 9. PREDICTION SET A ───────────────────────────────────────────────────────
\section{Prediction Set A — Jitter as Structural Information (Level I)}

\subsection{A1 — Jitter Discrimination Threshold}

\begin{predbox}{Template A1}

\textbf{Full title.}
Detection of temporal jitter in harmonic complexes matched for spectral envelope and pitch.

\subsubsection{Hypothesis}

Listeners can discriminate natural jitter (period variation $\approx 0.5$–$3\%$)
from perfectly periodic harmonics when spectral envelope, fundamental frequency,
and amplitude envelope are held constant.  This tests whether jitter constitutes
an audible signal dimension orthogonal to spectrum and amplitude.

\subsubsection{Falsifiable prediction}

\textbf{Primary (confirmatory):} Mean $d' > 1.0$ across participants (one-tailed
$t$-test against 1.0; $\alpha = 0.05$, power $= 0.80$).  A threshold of $d' = 1.0$
corresponds to approximately 76\% correct in 2AFC — a level of sensitivity that
requires jitter to be more than marginally detectable.  The lower value of
$d' > 0.5$ ($\approx 60\%$ correct) was rejected during design because it is
consistent with chance-level performance in many auditory discrimination studies
and would not demonstrate that jitter carries practically meaningful signal.

\textbf{Secondary (confirmatory):} $d'$ increases monotonically with jitter
magnitude across three pre-specified levels (0.5\%, 1.5\%, 2.5\% RMS period
variation).  This is tested via a one-way repeated-measures ANOVA with a
significant linear contrast ($p < 0.05$).  Both criteria must be met for the
prediction to be considered confirmed.

\subsubsection{Null hypothesis}

Mean $d' \leq 1.0$, or no monotonic relationship between $d'$ and jitter magnitude.

\subsubsection{Design}

\textbf{Stimuli.}
Synthesized harmonic complex ($f_0 = 220$~Hz, harmonics 1–10, equal amplitude
modulo a fixed spectral envelope, duration 500~ms, $\cos^2$ onset/offset ramps
of 10~ms).
\begin{itemize}
  \item \textbf{Condition A (three levels).}  Temporal jitter is generated
    algorithmically (see below) at RMS period variation of 0.5\%, 1.5\%, and
    2.5\% of mean period.  Applied by resynthesis from a period sequence; no
    recorded audio is used directly.
  \item \textbf{Condition B.}  Identical resynthesis procedure with all period
    deviations set to zero (perfectly isochronous).
\end{itemize}

\textbf{Jitter synthesis procedure.}
For target RMS jitter $J \in \{0.005, 0.015, 0.025\}$: generate $K = 111$
period values $T_k = T_0(1 + \epsilon_k)$, where $T_0 = 1/220$~s,
$\epsilon_k \sim \mathcal{N}(0, J^2)$, clipped to $[-3J, 3J]$.  Apply a
2nd-order Butterworth low-pass filter at 30~Hz to $\epsilon_k$ before use
(removes inaudibly fast variation while preserving temporal fluctuation in
the audible jitter range).  Resynthesize each period as one cycle of the
harmonic complex using randomized starting phases (per-file seed, fixed
across conditions).  Phase randomization controls waveform morphology and
fine temporal microstructure across conditions but does not equate spectral
envelopes, which are explicitly matched.  Stimuli are generated deterministically
from a fixed seed list; synthesis code is deposited to OSF prior to data collection.

\textbf{Period estimation (for replicator verification):}
Jitter magnitude in synthesized stimuli should be verified using an
autocorrelation-based or Hilbert instantaneous-phase period tracker,
not raw peak detection.  Peak detection in complex harmonic stimuli is
unreliable because harmonic interference displaces apparent peak locations.
The verification criterion is: measured $J$ within $\pm 0.002$ of target $J$.

\textbf{Task.}
Two-interval, two-alternative forced choice (2I-2AFC) with corrective feedback
after each trial.  Each interval is 500~ms, inter-stimulus gap 300~ms.
$X$ drawn uniformly from \{A, B\}; participant indicates which interval contained
jitter.

\textbf{Participants.}
$N = 30$ normal-hearing adults (audiometric threshold $\leq 20$~dBHL at 0.5–4~kHz).
Power analysis: $\Delta d' = 1.0$ within-subjects design, 90\% power implies
$n \approx 22$; 35\% buffer applied.

\textbf{Trials.}
120 per participant (30 per jitter level $\times$ 3 levels $+$ 30 jitter-absent),
randomized without replacement per block of 12.

\textbf{Blinding.}
The experimenter running sessions is blind to the jitter magnitude assigned to
each stimulus file.  Stimulus files are labelled by a numeric code assigned by
a separate researcher.  Analysis is conducted on anonymized data files before
the code is broken.

\subsubsection{Analysis}

Sensitivity is computed as:
\[
  d' = z(H) - z(F)
\]
where $H$ is the hit rate and $F$ is the false alarm rate.  Signal detection
theory is used because $d'$ is invariant to response bias, preventing
criterion differences between participants from inflating or suppressing the
sensitivity estimate.

\textbf{Primary:} One-sample $t$-test of mean $d'$ (collapsed across jitter
levels) against 1.0, one-tailed.

\textbf{Secondary:} One-way repeated-measures ANOVA on $d'$ at 0.5\%, 1.5\%,
and 2.5\% RMS jitter; pre-specified linear contrast.

\textbf{Exclusion criteria:} fewer than 80\% of trials completed; accuracy
$< 55\%$ on the 10-trial practice block (accuracy is used here rather than
$d'$ because the practice block is too short for stable SDT estimation).

\textbf{Confirmatory:} both primary and secondary criteria must be met.
\textbf{Exploratory:} interaction between jitter magnitude and musical training.

\end{predbox}

\subsection{A2 — Marine Coherence Predicts Auditory Salience}

\begin{predbox}{Template A2}

\textbf{Full title.}
Temporal coherence (Marine coherence metric $\kappa$) adds incremental predictive
power over spectral prominence in auditory salience ratings.

\subsubsection{Hypothesis}

The Marine coherence metric $\kappa$ predicts listener salience ratings beyond
what is accounted for by spectral prominence alone.

\subsubsection{Falsifiable prediction}

In a multiple regression, $\hat{\beta}_2 > 0$ with $p < 0.001$, and
$\Delta R^2 \geq 0.08$ after adding $\kappa$ to a model containing only
prominence $P$.  The threshold of $\Delta R^2 = 0.08$ reflects published
benchmarks in auditory salience and psychoacoustic modelling (e.g., Klapuri,
2006; Kayser et al., 2005) where a single novel predictor accounting for
8–10\% of unique variance is considered practically meaningful.  A higher
threshold ($\Delta R^2 \geq 0.15$) was considered but rejected as poorly
calibrated to the expected magnitude of a single jitter-derived metric relative
to the strong baseline provided by spectral prominence.

\subsubsection{Null hypothesis}

$\hat{\beta}_2 \leq 0$ or $\Delta R^2 < 0.04$ (coherence adds no practically
meaningful predictive power beyond spectral prominence).

\subsubsection{Design}

\textbf{Stimuli.}
100 partials extracted from 10 recorded instrumental notes (piano, violin, guitar,
voice; notes selected to span the full range of jitter magnitudes observed in
pilot measurements).  For each partial:
\begin{itemize}
  \item Prominence $P$ = dB above local noise floor (estimated via median filter
    of width 100~Hz in log-frequency).
  \item Marine coherence metric $\kappa = 1 / (J + 0.001)$, measured over a
    200~ms central window.
\end{itemize}

\textbf{Rating task.}
20 normal-hearing listeners rate each partial in isolation on a 0–100 continuous
slider labelled ``barely audible'' to ``very salient.''  Each partial is presented
alone (without harmonic context) once, in individually randomized order.  This
presentation mode is acknowledged as a limitation: partial salience in isolation
may differ from salience in a harmonic context.  A naturalistic follow-up study
is designated exploratory.  Mean rating across listeners is the dependent variable.

\textbf{Blinding.}
Listeners are blind to all hypotheses.  The analyst computing $\kappa$ is
separate from the analyst running the regression, to prevent confirmation bias
in feature extraction choices.

\textbf{Regression models.}
\begin{align}
  \text{Model 1:} \quad & \hat{S} = \beta_0 + \beta_1 P \\[4pt]
  \text{Model 2:} \quad & \hat{S} = \beta_0 + \beta_1 P + \beta_2 \kappa
\end{align}

The incremental test evaluates whether $\kappa$ explains variance in salience
over and above spectral prominence.  This is the appropriate test because the
paradigm claims that temporal coherence is an \emph{orthogonal} salience
dimension, not a substitute for spectral energy.

\subsubsection{Analysis}

Linear mixed model with listener as a random intercept, or ordinary least squares
on listener-averaged ratings.  Report $\Delta R^2$ and $F$-test for model
comparison.  Outlier removal: Cook's $D > 4/n$ for any partial, declared before
unblinding.

\textbf{Confirmatory:} $\hat{\beta}_2 > 0$, $p < 0.001$, and $\Delta R^2 \geq 0.08$.
\textbf{Exploratory:} interaction $P \times \kappa$; instrument-family moderator;
naturalistic harmonic-context replication.

\end{predbox}

\subsection{A2b — Harmonic-Context Replication (Secondary Confirmatory)}

\begin{predbox}{Template A2b}

\textbf{Full title.}
Marine coherence metric predicts salience of harmonic complexes presented in
full harmonic context.

\subsubsection{Rationale}

A2 presents partials in isolation, which is acknowledged as ecologically
limited: temporal coherence effects in auditory perception often require
harmonic context to manifest via common-fate grouping and fusion.  A2b is a
pre-registered secondary confirmatory study using naturalistic stimuli.  Both
A2 and A2b must be reported; they are independent tests of the same hypothesis
under different stimulus conditions.

\subsubsection{Falsifiable prediction (identical criterion to A2)}

$\hat{\beta}_2 > 0$, $p < 0.001$, and $\Delta R^2 \geq 0.08$.

\textbf{Additional criterion (demotes to exploratory if violated):}
The $\Delta R^2$ in A2b must be numerically $\geq$ the $\Delta R^2$ in A2.
If context adds no incremental effect, the harmonic-context result is
exploratory regardless of significance.

\subsubsection{Null hypothesis}

$\hat{\beta}_2 \leq 0$, or $\Delta R^2 < 0.08$, or $\Delta R^2_\text{A2b} <
\Delta R^2_\text{A2}$.

\subsubsection{Design}

\textbf{Stimuli.}  50 harmonic complexes (harmonics 1–10 simultaneously) from
the same 10 instrumental recordings used in A2.  Duration 1~s, 10~ms $\cos^2$
ramps.

\textbf{Feature aggregation.}  Two pre-registered strategies (both must meet
criteria for confirmation):
\begin{itemize}
  \item \textbf{Max rule:} $P_\text{cmplx} = \max_h P_h$;
    $\kappa_\text{cmplx} = \max_h \kappa_h$.
  \item \textbf{Energy-weighted mean:} $P_\text{cmplx} = \sum_h w_h P_h$;
    $\kappa_\text{cmplx} = \sum_h w_h \kappa_h$, where
    $w_h = P_h / \sum_h P_h$.
\end{itemize}

\textbf{Rating task.}  20 normal-hearing listeners (may overlap with A2 sample).
0–100 slider; mean across listeners is the dependent variable.

\textbf{Analysis.}  Same mixed model as A2; report $\Delta R^2$ per aggregation
rule.  Confirmation requires both aggregation rules to meet all three criteria.

\textbf{Confirmatory:} as stated.
\textbf{Exploratory:} comparison of max vs.\ weighted-mean $\Delta R^2$;
partial vs.\ complex moderation.

\end{predbox}

% ── 10. PREDICTION SET B ──────────────────────────────────────────────────────
\section{Prediction Set B — HCF1 Compression Architecture (Level II)}

\subsection{B1 — Harmonic Envelope Conformity}

\begin{predbox}{Template B1}

\textbf{Full title.}
Amplitude envelopes of upper harmonics are linearly predictable from the
fundamental in acoustic instruments.

\subsubsection{Hypothesis}

For physically coupled harmonic oscillators, the amplitude envelope of each upper
harmonic is well approximated by a linear function of the fundamental's envelope.
This is the necessary prior condition for the \hcf{} compression scheme.

\subsubsection{Falsifiable prediction}

Median $R^2 > 0.90$ across harmonics $h = 2, \ldots, 10$, with the lower bound
of the 95\% bootstrap confidence interval exceeding 0.80.

\subsubsection{Null hypothesis}

Median $R^2 < 0.70$ (harmonic-specific variation dominates; envelopes are not
jointly predictable from the fundamental).

\subsubsection{Design}

\textbf{Recordings.}
20 notes per instrument family (different pitches, sampled at 44.1~kHz, minimum
duration 2~s): plucked string (classical guitar), struck string (piano), bowed
string (violin), struck membrane (timpani).

\textbf{Analysis procedure.}
\begin{enumerate}
  \item Detect $f_0$ per frame (10~ms hop, 50\% overlap) via an autocorrelation
    method.
  \item Extract amplitude envelope $e_h(t)$ for harmonics $h = 1, \ldots, 10$
    via an 8th-order Butterworth bandpass filter (bandwidth $\pm 20\%$ of $hf_0$),
    followed by full-wave rectification and low-pass filtering at 30~Hz.
  \item For each harmonic $h = 2, \ldots, 10$, fit by ordinary least squares:
    \[
      \hat{e}_h(t) = a_h \cdot e_1(t) + b_h
    \]
    Compute $R^2 = 1 - \mathrm{SS}_\text{res} / \mathrm{SS}_\text{tot}$.
\end{enumerate}

\textbf{Exclusion:} Notes with inharmonicity $> 1\%$ are excluded from the
confirmatory sample and reported in a separate table.  Inharmonicity is
operationally defined as the mean absolute deviation of measured partial
frequencies from their predicted values $h \cdot f_0$, normalized by $f_0$:
\[
  \mathrm{Inh} = \frac{1}{H} \sum_{h=1}^{H}
    \frac{|\hat{f}_h - h \cdot f_0|}{f_0}
\]
where $\hat{f}_h$ is the measured frequency of the $h$-th partial (estimated
via spectral peak picking on a zero-padded FFT).  This definition is fixed
prior to data collection; the threshold of 1\% is pre-registered and cannot
be adjusted post-hoc.

A \textbf{sensitivity analysis} is mandatory: the confirmatory $R^2$ test is
re-run including all excluded notes.  If the result changes sign or crosses the
threshold boundary, this must be reported prominently.

\subsubsection{Analysis}

Median $R^2$ and 95\% bootstrap confidence interval (10\,000 resamples over
individual recordings).  Secondary: distribution of slope $a_h$ — expected to
decrease monotonically with $h$ for most instrument families.

\textbf{Adaptive filter bandwidth (pre-registered exception):}
If a partial's instantaneous $f_0$ deviates $> 15\%$ from $h \cdot f_0$ for
$> 10\%$ of frames (e.g., notes with wide vibrato or pitch drift), the
Butterworth filter bandwidth is widened to $\pm 30\%$ of $h f_0$ for that
partial.  The proportion of frames triggering this widening must be reported
per note.  This exception is declared before analysis; post-hoc bandwidth
adjustment is not permitted.

\textbf{Confirmatory:} $R^2$ threshold.
\textbf{Exploratory:} differences across instrument families; interaction with
note duration and pitch register.

\end{predbox}

\subsection{B2 — HCF1 Bit Reduction and Perceptual Transparency}

\begin{predbox}{Template B2}

\textbf{Full title.}
\hcf{} with 1-bit conformity flags achieves $\geq\!60\%$ bit reduction without
audible degradation in a blind listening test.

\subsubsection{Hypothesis}

Storing one master envelope and a 1-bit conformity flag per harmonic yields
substantial bitrate savings while remaining perceptually indistinguishable from
the original.

\subsubsection{Falsifiable prediction}

(1) Bit reduction $\geq 60\%$, \textbf{and} (2) ABX discrimination accuracy does
not exceed chance: binomial test $p > 0.05$; equivalently, proportion correct
$< 0.575$ (upper 95\% one-sided CI for $n = 50$, $p_0 = 0.5$).

\subsubsection{Null hypotheses}

(1) Bit reduction $< 40\%$, \textbf{or} (2) ABX accuracy $> 57.5\%$.

\subsubsection{Design}

\textbf{Compression conditions.}

\emph{Baseline.}  Store envelopes $e_1, \ldots, e_{10}$ independently at 16~bits,
100~Hz update rate, yielding:
\[
  R_\text{baseline} = 10 \times 16 \times 100 = 16{,}000~\text{bits/s}
\]

\emph{\hcf{}.}  Store $e_1$ at 16~bits, 100~Hz.  For $h = 2, \ldots, 10$, store
one bit per frame:
\[
  \text{flag}_{h,t} = \begin{cases}
    0 & \text{(conform)} \Rightarrow \hat{e}_h(t) = a_h e_1(t) + b_h \\
    1 & \text{(nonconform)} \Rightarrow \text{store residual at 16 bits}
  \end{cases}
\]

Effective bitrate:
\[
  R_\text{HCF1} = \underbrace{16 \times 100}_{\text{fundamental}} +
                  \underbrace{9 \times 1 \times 100}_{\text{9 flag bits/s}} +
                  f_\text{nc} \times 9 \times \underbrace{16 \times 100}_{\text{residual/harmonic}}
                  ~\text{bits/s}
\]
where $f_\text{nc} \in [0,1]$ is the empirically measured nonconformity
fraction.  At zero nonconformity this reduces to $1600 + 900 = 2500$~bits/s;
at full nonconformity it approaches $1600 + 900 + 14400 = 16900$~bits/s
($\approx$~baseline).  The 60\% reduction criterion requires that the
empirically measured $f_\text{nc}$ yield $R_\text{HCF1} \leq 0.4 \times
16000 = 6400$~bits/s.

\textbf{Listening test.}
10 audio files (5–15~s each): solo speech, solo piano, jazz trio, string quartet,
pop vocal, acoustic guitar, orchestral excerpt, electronic music, percussion
ensemble, mixed ensemble.  20 normal-hearing listeners; 50 ABX trials each.
$A = \text{original}$, $B = \text{\hcf{}-reconstructed}$, $X = $ random draw.

\subsubsection{Analysis}

Bit reduction computed directly from measured bitrates.  ABX: binomial test
against $p_0 = 0.5$; non-inferiority margin $= 0.075$.  Pass criterion: reduction
$\geq 60\%$ \textbf{and} upper CI of proportion correct $< 57.5\%$.

\textbf{Confirmatory:} both criteria jointly.
\textbf{Exploratory:} effect of musical genre on reduction ratio; correlation
between $R^2$ from B1 and \hcf{} efficiency per recording.

\end{predbox}

% ── 11. PREDICTION SET C ──────────────────────────────────────────────────────
\section{Prediction Set C — \mem{} Associative Memory Architecture (Level III)}

\subsection{C1 — Quadratic Storage Capacity via Phase Encoding}

\begin{predbox}{Template C1}

\textbf{Full title.}
Sharp retrieval transition in \mem{}: testing predicted dense relational
capacity in a phase-based associative memory architecture.

\subsubsection{Hypothesis}

A \mem{} system with $N$ wave functions and a pairwise offset matrix
$\Delta_{ij}$ is hypothesised to support dense relational encoding such that
retrieval accuracy remains high below approximately $M^2/N \approx 1$ and
degrades sharply near $M^2/N \approx 1$--$2$.  This is motivated by analogy
to constraint satisfaction phase transitions (e.g., random $k$-SAT), where
underconstrained systems admit solutions near-perfectly and overconstrained
systems fail sharply.  The hypothesis is not that a formal proof of $O(N^2)$
capacity exists, but that the empirical accuracy curve behaves consistently
with this scaling regime in simulations up to $N = 256$.

\subsubsection{Falsifiable predictions}

\begin{itemize}
  \item Retrieval accuracy $> 80\%$ when $M^2 \leq N$ (underconstrained regime).
  \item Retrieval accuracy $< 50\%$ when $M^2 > 2N$ (overconstrained regime).
  \item \textbf{Sharpness (pre-registered quantitative criterion):}  A logistic
    function $\hat{a}(x) = 1/(1 + e^{-k(x - x_0)})$ is fit to accuracy as a
    function of $x = M^2/N$.  The transition is considered sharp if the fitted
    slope parameter $k > 8$ (relaxed from 10 to account for finite-size effects
    at $N = 64$).  Alternatively, if the logistic fit converges poorly
    ($R^2_\text{fit} < 0.80$), the bootstrapped accuracy difference
    $a(M^2 = N) - a(M^2 = 2N) > 0.4$ is used.  The primary criterion
    (logistic slope) takes precedence; the fallback is declared before analysis.
\end{itemize}

These thresholds are empirical predictions, not mathematical guarantees.
Random constraint systems can become frustrated before naive counting thresholds
suggest, and finite-size effects may shift the observed transition.  The
predictions specify what the simulations \emph{should} show if the architecture
behaves as hypothesised; deviations are informative regardless of direction.

\subsubsection{Null hypothesis}

Accuracy scales linearly with $N$ (no quadratic advantage), or $k \leq 5$ in the
logistic fit (gradual degradation rather than sharp phase transition).

\subsubsection{Design}

\textbf{Implementation.}
Floating-point (double precision) simulation of \mem{} with
$N \in \{64, 128, 256\}$.  The architecture uses two components:
(1) $N$ wave functions $\psi_k = e^{i\phi_k}$ on the unit circle, and
(2) an antisymmetric offset matrix $\Delta \in \mathbb{R}^{N \times N}$
(zero diagonal) storing pairwise target offsets.  It is the offset matrix,
not the oscillator states alone, that encodes associations.

\medskip
\noindent\textbf{\mem{} pseudocode (corrected, replication-ready):}

\smallskip
\noindent\textit{Initialization:}
\begin{itemize}[topsep=1pt, itemsep=0pt]
  \item $\phi_k \sim \mathrm{Uniform}[0, 2\pi)$ for $k = 1, \ldots, N$.
  \item $\Delta_{ij} \leftarrow 0$ for all $i, j$; $\Delta_{ii} = 0$ (enforced).
\end{itemize}

\noindent\textit{Storage of association $(i, j)$:}
\begin{enumerate}[topsep=2pt, itemsep=1pt]
  \item Record target offset: $\Delta_{ij} \leftarrow \angle(\psi_j / \psi_i)$;
    $\Delta_{ji} \leftarrow -\Delta_{ij}$ (antisymmetric).
  \item Run relaxation (10 iterations, fixed):
    \begin{enumerate}[label=\alph*., topsep=1pt, itemsep=0pt]
      \item Compute current pairwise angles:
        $C_{ij} = \angle(\psi_j / \psi_i)$ for all $(i,j)$.
      \item Error matrix: $E_{ij} = \angle(e^{i(\Delta_{ij} - C_{ij})})$
        (wrapped to $[-\pi, \pi]$).
      \item Update each wave function:
        $\phi_k \mathrel{+}= \eta \cdot \mathrm{clip}
          \!\left(\tfrac{1}{N}\sum_{j}E_{kj} + \tfrac{1}{N}\sum_{i}E_{ik},\;
          {-0.5},\; {0.5}\right)$
        with learning rate $\eta = 0.1$.
      \item Normalize: $\psi_k \leftarrow \psi_k / |\psi_k|$.
    \end{enumerate}
\end{enumerate}

\noindent\textit{Retrieval given cue $i$:}
\begin{enumerate}[topsep=2pt, itemsep=1pt]
  \item For each candidate $j \neq i$, compute predicted phase:
    $\hat{\phi}_j = \phi_i + \Delta_{ij}$.
  \item Angular distance: $d_j = |\angle(e^{i(\hat{\phi}_j - \phi_j)})|$.
  \item Retrieved item: $j^* = \arg\min_{j \neq i} d_j$.
  \item Correct if $j^* = j_\text{true}$.
\end{enumerate}

This specification is complete and self-contained.  Any implementation that
deviates from this pseudocode — including different relaxation iterations,
learning rate schedules, or distance metrics — constitutes a variant and must
be so described.

\textbf{Load conditions.}
$M \in \bigl\{
  \lfloor\sqrt{N/2}\rfloor,
  \lfloor\sqrt{N}\rfloor,
  \lfloor\sqrt{2N}\rfloor,
  \lfloor\sqrt{4N}\rfloor
\bigr\}$,
yielding four load regimes per $N$.  Association pairs are drawn uniformly at
random from the $N$ item slots without replacement within each set.

\textbf{Trials.}
100 independently generated random association sets per $(N, M)$ condition.

\subsubsection{Analysis}

Accuracy plotted as a function of $M^2/N$ with 95\% bootstrap CI over the
100 association sets.  Logistic regression fit to the accuracy curve;
slope $k$ extracted and tested against the pre-registered threshold of 8.

\textbf{Validation test (mandatory before full C1 run):}  Run $N = 16$,
$M^2/N \in \{0.25, 0.5, 1.0, 1.5, 2.0, 3.0\}$, 20 seeds each.  Plot
accuracy vs.\ $M^2/N$.  If no transition is visible (accuracy remains
constant $> 80\%$ or $< 50\%$ throughout), the implementation is incorrect
and must be debugged before proceeding.

\textbf{Confirmatory:} accuracy thresholds at $M^2 \approx N$ and
$M^2 \approx 2N$; logistic sharpness criterion $k > 8$.
\textbf{Exploratory:} effect of phase quantisation (fixed-point vs.\
floating-point); comparison to complex-valued Hopfield networks (Noest, 1988;
Jankowski et al., 1996) on the same task.

\end{predbox}

\subsection{C2 — Resistance to Catastrophic Forgetting}

\begin{predbox}{Template C2}

\textbf{Full title.}
\mem{} retains prior associations after sequential learning, unlike Hopfield
networks at equivalent load.

\subsubsection{Hypothesis}

\mem{} shows $\leq 10\%$ drop in recall accuracy after learning a second
association set; a standard Hopfield network at the same load drops $\geq 50\%$.
This tests whether the phase-encoding strategy confers a stability-plasticity
advantage.

\subsubsection{Falsifiable predictions}

\begin{itemize}
  \item \mem{} retention $\geq 90\%$ of original accuracy.
  \item Hopfield retention $\leq 50\%$ of original accuracy.
  \item Difference significant at $p < 0.001$ (two-sample Welch $t$-test).
\end{itemize}

\subsubsection{Null hypothesis}

\mem{} retention $< 80\%$, or not statistically different from Hopfield at
$\alpha = 0.001$.

\subsubsection{Design}

\textbf{Networks.}
\mem{} ($N = 256$ complex wave functions, retrieval as specified in C1) vs.\
a \textbf{pair-association Hopfield baseline} ($N = 256$ binary units).

The Hopfield baseline is implemented as a pair-association task, not a pattern
completion task, to ensure the comparison is fair.  Concretely: each item $i$
is represented as a random $N/2$-bit binary vector $\mathbf{v}_i$; each pair
$(i,j)$ is stored as the outer-product Hebbian update on the weight matrix
$W \mathrel{+}= \mathbf{v}_i \mathbf{v}_j^T + \mathbf{v}_j \mathbf{v}_i^T$.
Retrieval: given cue $\mathbf{v}_i$, one step of synchronous update yields
$\hat{\mathbf{v}}_j = \mathrm{sgn}(W \mathbf{v}_i)$; correct if nearest to
$\mathbf{v}_j$.  This formulation is equivalent to standard Hopfield pattern
completion on the concatenated pair representation but expressed in the
association-retrieval frame.

Both architectures are therefore solving the \emph{same task} under the same
performance metric (proportion of correctly retrieved partners), differing only
in their internal representations.

\textbf{Procedure.}
\begin{enumerate}
  \item Learn association set $A$ ($M = 16$ pairs); measure accuracy on $A$.
  \item Learn association set $B$ ($M = 16$ new, non-overlapping pairs, using
    the same update rule without resetting the weight matrix / phase registers).
  \item Retest accuracy on $A$ (same cues as step 1).
\end{enumerate}

\[
  \text{Retention} = \frac{\text{accuracy on } A \text{ after step 3}}
                          {\text{accuracy on } A \text{ after step 1}} \times 100\%
\]

\textbf{Repetitions.}  50 independent random seeds per network type (different
random item vectors / phase initializations).

\subsubsection{Analysis}

Two-sample Welch $t$-test comparing \mem{} vs.\ Hopfield retention distributions.
Also report: proportion of seeds where \mem{} exceeds Hopfield by $> 20\%$.
Scatter plot of original vs.\ retained accuracy (one point per seed, per network).

\textbf{Confirmatory:} retention thresholds and network comparison.
\textbf{Exploratory:} retention as a function of increasing load
($M \in \{32, 64\}$); comparison to other continual learning baselines
(e.g., Elastic Weight Consolidation).

\end{predbox}

% ── 12. INTERPRETATION OF NEGATIVE RESULTS ────────────────────────────────────
\section{Interpretation of Negative Results}
\label{sec:negative}

A pre-registration that does not specify the consequences of falsification is
incomplete.  The following table states what each negative outcome would and would
not entail.

\noindent\begin{minipage}{\textwidth}
\centering
\renewcommand{\arraystretch}{1.5}
\begin{tabular}{p{1.1cm} p{4.2cm} p{5.8cm}}
\toprule
\textbf{Pred.} & \textbf{If falsified (negative outcome)} & \textbf{Implications and limits} \\
\midrule
A1 &
  Listeners cannot detect jitter at $d' > 1.0$, or $d'$ does not increase
  monotonically with jitter magnitude &
  Weakens the claim that jitter encodes perceptually salient structural
  information at practically meaningful levels.  Note: $d' > 0.5$ (weak
  detection) might still occur; the higher threshold is deliberately set to
  avoid claiming support from marginal sensitivity.  Does not falsify B or C. \\
A2 &
  $\kappa$ adds no salience variance ($\Delta R^2 < 0.08$, $p \geq 0.001$) &
  Weakens the Marine coherence metric as a perceptual predictor.  Does not
  falsify the compression claim (B) or memory claim (C).  Would suggest
  $\kappa$ is an engineering quantity without direct perceptual analogue,
  but does not preclude its use as a compression or detection feature. \\
B1 &
  Harmonic envelopes are not linearly predictable ($R^2 < 0.70$) &
  Directly undermines the conformity assumption behind \hcf{}.  Makes B2
  theoretically unmotivated.  Does not affect Level~I (A) or Level~III (C)
  predictions. \\
B2 &
  \hcf{} achieves $< 40\%$ reduction or is audibly distinguishable &
  Falsifies the specific engineering claim; does not falsify B1.  May indicate
  that the conformity fraction is instrument- or genre-specific rather than
  general. \\
C1 &
  No quadratic capacity advantage; logistic slope $k \leq 5$; or accuracy does
  not drop below 50\% at $M^2 = 2N$ &
  Weakens the central computational claim of \mem{}.  Does not affect Level~I
  or Level~II predictions.  Even if C1 fails, the geometric interpretability
  of phase-based encoding (Section~4) remains a valid architectural property. \\
C2 &
  \mem{} shows comparable forgetting to Hopfield &
  Weakens the stability-plasticity claim.  Combined with C1 failure, would
  substantially undermine the Level~III architecture as specified.  Would not
  affect psychoacoustic or compression predictions. \\
\bottomrule
\end{tabular}
\captionof{table}{Predicted consequences of falsification for each pre-registered prediction.
  Predictions are empirically independent; falsification of one does not entail
  falsification of others.}
\end{minipage}

\begin{notebox}
\textbf{Independence of levels.}  The three prediction sets are architecturally
related (Section~3) but empirically independent.  A complete failure of Set~C
would leave Sets~A and~B entirely untouched, and vice versa.  This modularity
is a deliberate feature: it means the package generates genuine evidence regardless
of which predictions succeed or fail.
\end{notebox}

% ── 13. REPLICATION PATHWAYS ──────────────────────────────────────────────────
\section{Replication Pathways}

\subsection{Accessibility by prediction}

\noindent\begin{minipage}{\textwidth}
\centering
\renewcommand{\arraystretch}{1.4}
\begin{tabular}{p{1.0cm} p{2.8cm} p{3.5cm} p{4.2cm}}
\toprule
\textbf{Pred.} & \textbf{Equipment needed} & \textbf{Software} & \textbf{Est.\ runtime} \\
\midrule
A1  & Audiometric booth; headphones; standard PC &
      Python (NumPy, SciPy, PsychoPy) or MATLAB &
      $\approx 3$~h data collection; 30~min analysis \\
A2  & Recordings of acoustic instruments; standard PC &
      Python (librosa, statsmodels) or R &
      $\approx 1$~h data collection; 1~h analysis \\
B1  & Recordings (publicly available datasets acceptable) &
      Python (librosa, SciPy) or MATLAB &
      $\approx 2$~h analysis per instrument family \\
B2  & Audiometric booth; headphones; standard PC &
      Python (SciPy, soundfile, PsychoPy) &
      $\approx 3$~h data collection; 1~h analysis \\
C1  & Standard PC (no GPU required) &
      Python (NumPy, SciPy) &
      $\approx 2$~h total (all $N$ and $M$ conditions) \\
C2  & Standard PC (no GPU required) &
      Python (NumPy) &
      $\approx 30$~min total \\
\bottomrule
\end{tabular}
\captionof{table}{Replication requirements by prediction.  C1 and C2 are purely
  computational and require no human participants or specialist equipment.}
\end{minipage}
\subsection{Recommended entry points for independent labs}

Predictions C1 and C2 are the easiest entry points for independent replication:
they require only a standard Python environment, no human participants, no IRB
approval, and have runtime measured in minutes to hours.  Predictions B1 is the
next most accessible, as it requires only existing recordings (freely available
datasets such as the University of Iowa Musical Instrument Samples library are
suitable).  Predictions A1, A2, and B2 require a listening environment; a quiet
room with calibrated headphones is sufficient for A1 and A2, though an
audiometric booth is preferred.

\subsection{Recommended datasets}

\begin{itemize}
  \item \textbf{B1, A2 (instrumental recordings):}  University of Iowa Electronic
    Music Studios Musical Instrument Samples; NSynth (Engel et al., 2017);
    RWC Music Database (Goto et al., 2003).
  \item \textbf{B2 (audio quality):}  EBU SQAM test sequences; internally generated
    stimuli following the compression procedure specified above.
  \item \textbf{A1 (jitter extraction):}  Any high-quality close-microphone
    recording of an acoustic guitar; jitter parameters specified in Template A1.
\end{itemize}

% ── 14. SUMMARY TABLE ─────────────────────────────────────────────────────────
\section{Summary of Pre-Registered Predictions}

\noindent\begin{minipage}{\textwidth}
\centering
\renewcommand{\arraystretch}{1.5}
\begin{tabular}{p{1.1cm} p{3.5cm} p{3.5cm} p{2.0cm} p{2.2cm}}
\toprule
\textbf{Pred.} & \textbf{Dependent variable} & \textbf{Statistical criterion}
  & \textbf{Sample} & \textbf{Suggested venue} \\
\midrule
A1 & Sensitivity index $d'$ + monotonic magnitude effect &
  $d' > 1.0$; linear contrast $p < 0.05$ &
  30 listeners &
  JASA; Acta Acust. \\
A2 & $\Delta R^2$, $\hat{\beta}_2$ (isolated partials) &
  $\Delta R^2 \geq 0.08$, $p < 0.001$ &
  100 partials, 20 listeners &
  J.\ Acoust.\ Soc. \\
A2b & $\Delta R^2$, $\hat{\beta}_2$ (harmonic context) &
  $\Delta R^2 \geq 0.08$, both aggregation rules &
  50 complexes, 20 listeners &
  J.\ Acoust.\ Soc. \\
B1 & Median $R^2$ (harmonic conformity) &
  $R^2 > 0.90$ (95\% CI $> 0.80$) &
  80 recordings &
  ICASSP; EURASIP \\
B2 & Bit reduction + ABX accuracy &
  $\geq 60\%$ reduction \& correct $< 57.5\%$ &
  10 files, 20 listeners &
  AES; IEEE TASLP \\
C1 & Retrieval accuracy vs.\ load &
  $> 80\%$ at $M^2 \leq N$; $< 50\%$ at $M^2 > 2N$; slope $k > 8$ &
  100 sets per $(N,M)$ &
  Neural Comp.; NeurIPS \\
C2 & Retention (\%) after sequential learning &
  \mem{} $\geq 90\%$; Hopfield (corrected) $\leq 50\%$; $p < 0.001$ &
  50 seeds per network &
  Neural Comp.; ICLR \\
\bottomrule
\end{tabular}
\captionof{table}{Summary of pre-registered predictions with current statistical
  criteria, sample sizes, and suggested venues.  All thresholds reflect
  revisions made prior to data collection; the version history is recorded
  in the OSF repository.}
\end{minipage}

\subsection*{Shared methodological properties}

All six predictions satisfy the following requirements for confirmatory
pre-registration:

\begin{itemize}
  \item \textbf{Falsifiable.}  An explicit null hypothesis is stated for each,
    identifying what negative evidence would look like.
  \item \textbf{Operationalized.}  Key variables are defined with explicit
    measurement procedures: $d'$, $\kappa$, $R^2$, bit reduction, retrieval
    accuracy, retention.
  \item \textbf{Scope-limited.}  No prediction appeals to ``consciousness,''
    ``authenticity,'' or ``cognitive field'' — constructs that remain
    unoperationalized at Level~IV.
  \item \textbf{Reproducible.}  Each study is executable by an independent
    laboratory without privileged access to proprietary instruments or data.
  \item \textbf{Modular.}  The prediction sets are empirically independent;
    see Section~\ref{sec:negative} for a systematic account of what each
    negative outcome would and would not entail.
\end{itemize}

% ── 15. STANDARD BOILERPLATE ──────────────────────────────────────────────────
\section{Standard Pre-Registration Boilerplate}

\subsection{IRB and ethics}

Studies involving human participants (A1, A2, B2) require approval from an
Institutional Review Board or equivalent ethics committee before data collection.
Participants must provide written informed consent and be compensated at a rate
consistent with local norms.  Auditory stimuli should be presented at
$\leq 70$~dBSPL for all listening experiments to avoid risk of hearing damage.
Computational studies (B1, C1, C2) do not involve human participants and do not
require IRB review.

\subsection{Data availability statement}

Upon completion of each study, the following materials will be deposited to an
open repository (OSF, Zenodo, or equivalent under a CC-BY licence):
\begin{itemize}
  \item All stimuli, or synthesis code sufficient to reproduce them exactly.
  \item Raw participant data, anonymized in compliance with applicable regulations.
  \item Analysis scripts (Python, R, or MATLAB, with inline comments).
  \item This pre-registration document, version-controlled with commit hash.
\end{itemize}

\subsection{Deviations from the pre-registered plan}

Any deviation from the analysis pipeline specified above must be declared
explicitly in the manuscript, with a statement distinguishing confirmatory
(pre-registered) from exploratory (post-hoc) findings.  Exploratory analyses
are encouraged and should be reported transparently; the distinction is not a
judgment of scientific value but a statement of inferential status.

\subsection{Conflicts of interest}

The author has a prior theoretical investment in the \scp{}.  Independent
replication is therefore especially valuable, and replication teams are encouraged
to file their own pre-registration citing this document as the source of the
original predictions.

\subsection{Suggested repositories}

Pre-registrations may be filed at:
\begin{itemize}
  \item \href{https://osf.io}{Open Science Framework (OSF)} — recommended for the
    full package, as it supports version control and time-stamped registration.
  \item \href{https://aspredicted.org}{AsPredicted} — suitable for individual
    templates (A1–C2) filed separately.
  \item Journal-specific Registered Reports tracks at \emph{JASA},
    \emph{Psychological Science}, \emph{Neural Computation}, or \emph{ICLR}.
\end{itemize}

\vfill
{\color{navy}\rule{\linewidth}{0.8pt}}
{\centering\small\color{slategray}
Pre-Registration Package — Salience-Centered Paradigm (Levels~I–III)
\quad $\cdot$ \quad Flyxion (independent researcher)
\quad $\cdot$ \quad May 2026\par}

\end{document}
 
